\chapter{LEAN\_Warmup: proved tools and computations}

These modules contain completed mathematics as well as one unfinished mixed-derivative
formula. They remain in the exploratory library, with their current names and dependencies.

\section{Specialization and the factorization criterion}

\begin{lemma}[Specialization API]

\label{lem:setzero-api}

\lean{setZero_X, setZero_X_self, setZero_X_of_ne, setZero_C, setZero_monomial, support_setZero_subset, degreeOf_setZero_le, setZero_mem_supported_of_mem, setZero_eq_self_iff, setZero_setZero, setZero_comm, pderiv_setZero_self, pderiv_setZero_of_ne, setZero_map, eval_setZero}

\leanok

\uses{def:setzero,lem:affine-derivative}

Specializing $x_k$ to zero fixes constants and the other variables, kills $x_k$,
and kills exactly those monomials with a positive $k$th exponent. It decreases support
and individual degrees, preserves support restrictions, and fixes exactly the polynomials
independent of $x_k$. Specializations are idempotent and commute with one another.
Differentiation in $x_k$ after specialization is zero; differentiation in another variable
commutes with specialization. Specialization also commutes with coefficient maps and
has the expected evaluation formula.

\end{lemma}

\begin{proof}

\leanok

The linked Lean declarations contain the proof.

\end{proof}

\begin{lemma}[Two-variable affine decomposition]

\label{lem:bilinear}

\lean{exists_eq_X_mul_X_add, rayleighDiff_X_mul_X_add, rayleighDiff_mul_eq_zero}

\leanok

\uses{lem:affine-decompose,lem:setzero-api}

If $f$ has degree at most one in $x_i$ and $x_j$, it can be written
\[ f=ax_ix_j+bx_i+cx_j+d \]
with $a,b,c,d$ supported outside $\{i,j\}$. For $i\ne j$, the Rayleigh difference of
this expression is $bc-ad$. If $f=gh$ with $g$ independent of $x_i$ and $h$ independent
of $x_j$, then $\Delta_{ij}(f)=0$, without degree hypotheses.

\end{lemma}

\begin{proof}

\leanok

The linked Lean declarations contain the proof.

\end{proof}

\begin{lemma}[Rayleigh vanishing and factorization (Lemma 2.2)]

\label{lem:22}

\lean{rayleighDiff_eq_zero_iff}

\leanok

\uses{lem:bilinear,lem:support}

Over a unique factorization domain, if $f$ is affine in $x_i$ and $x_j$, then
\[ \Delta_{ij}(f)=0\quad\Longleftrightarrow\quad
 f=gh,\quad g\in R[x_k:k\ne i],\quad h\in R[x_k:k\ne j]. \]
The formal statement also allows $i=j$.

\end{lemma}

\begin{proof}

\leanok

For distinct indices use $bc=ad$. If $a\ne0$, unique factorization splits
$a=\gamma\alpha$, $b=\gamma\beta$, $c=\alpha\delta$, so $d=\beta\delta$ and
$f=(\alpha x_j+\beta)(\gamma x_i+\delta)$. Divisor-support lemmas give the support
restrictions. If $a=0$, one of $b,c$ vanishes. For equal indices, the square of the
derivative vanishes in a domain, and affine decomposition finishes the proof.

\end{proof}

\section{Differential operators and mixed derivatives}

\begin{lemma}[Squarefree exponent API]

\label{lem:exponent-api}

\lean{Finset.exponent_empty, Finset.exponent_insert, Finset.exponent_apply}

\leanok

\uses{def:multiaffine}

The exponent vector $e_S$ is zero on the empty set, adding a new element adds its
single-coordinate exponent, and $e_S(i)$ is one or zero according as $i\in S$ or not.

\end{lemma}

\begin{proof}

\leanok

The linked Lean declarations contain the proof.

\end{proof}

\begin{definition}[Evaluation in a noncommutative algebra]

\label{def:commuting-eval}

\lean{MvPolynomial.aevalOfCommute, MvPolynomial.aevalOfCommute_X}

\leanok

A family of pairwise commuting elements in an algebra defines an algebra homomorphism
from the multivariate polynomial algebra, even if the target algebra itself is not
commutative. Variables evaluate to the supplied elements.

\end{definition}

\begin{definition}[Polynomial differential operators]

\label{def:differential-op}

\lean{MvPolynomial.pderivOp, MvPolynomial.pderivOp_X, MvPolynomial.pderivOp_mul_apply, MvPolynomial.pderivOp_X_pow}

\leanok

\uses{def:commuting-eval}

Substituting the commuting partial derivatives for variables defines
$p\mapsto p(\partial)$ with values in linear endomorphisms of the polynomial ring.
A variable gives its partial derivative; polynomial multiplication gives composition;
$x_i^k$ gives the $k$th iterate of $\partial_i$.

\end{definition}

\begin{lemma}[Coefficients of differential operators]

\label{lem:operator-coeff}

\lean{MvPolynomial.coeff_pderiv_iterate, MvPolynomial.coeff_pderivOp_monomial, MvPolynomial.constantCoeff_pderivOp_monomial, MvPolynomial.coeff_pderivOp_prod_X, MvPolynomial.constantCoeff_pderivOp_prod_X, MvPolynomial.constantCoeff_pderivOp_apply_monomial, MvPolynomial.pderivOp_injective}

\leanok

\uses{def:differential-op,lem:exponent-api}

The coefficient of $x^t$ in $\partial_i^k f$ is
$(t_i+k)^{\underline{k}}[x^{t+ke_i}]f$. For a multi-index $s$, the coefficient in
$\partial^s f$ is $\prod_i(t_i+s_i)^{\underline{s_i}}[x^{t+s}]f$.
At $t=0$ this is $\prod_i s_i!\,[x^s]f$.
For squarefree $s=e_S$ no factorial remains at zero, and
$[x^t]\partial_S f=\prod_{i\in S}(t_i+1)[x^{t+e_S}]f$.
Conversely, the constant coefficient of $p(\partial)x^s$ is $\prod_i s_i!\,[x^s]p$.
Over a characteristic-zero domain these formulas make $p\mapsto p(\partial)$ injective.

\end{lemma}

\begin{proof}

\leanok

The linked Lean declarations contain the proof.

\end{proof}

\begin{lemma}[Mixed derivatives of multiaffine polynomials]

\label{lem:mixed-support}

\lean{coeff_add_exponent_ne_zero_of_coeff_pderivOp_ne_zero, pderivOp_prod_X_mem_supported, IsMultiaffine.pderivOp_prod_X}

\leanok

\uses{lem:operator-coeff,def:multiaffine}

A nonzero coefficient of $\partial_S f$ forces the corresponding shifted coefficient
of $f$ to be nonzero. For multiaffine $f$, $\partial_S f$ is supported outside $S$
and remains multiaffine.

\end{lemma}

\begin{proof}

\leanok

The linked Lean declarations contain the proof.

\end{proof}

\begin{lemma}[Generalized Jacobi formula]

\label{lem:mixed-det}

\lean{Matrix.pderivOp_detPoly}

\leanok

\uses{def:differential-op,lem:detpoly-derivatives}

The proposed identity is
\[ \partial_S f_A=\det M_A[S^c,S^c]. \]
The proof draft inducts on $S$; the empty and insertion steps remain unfinished.
The first- and second-derivative library formulas are completed special cases.

\end{lemma}

\noindent\emph{Status: unfinished proof.}

\begin{lemma}[Constant coefficient of a polynomial principal minor]

\label{lem:constant-minor}

\lean{Matrix.constantCoeff_principalMinor_mvCharmatrix}

\leanok

\uses{def:fiber,def:detpoly}

The constant coefficient of $\det M_A[S,S]$ is $A_S$.

\end{lemma}

\begin{proof}

\leanok

The linked Lean declarations contain the proof.

\end{proof}



\section{Computable multivariate polynomials}

\begin{definition}[A computation model for polynomials]

\label{def:computable}

\lean{CMvPolynomial, CMvPolynomial.X, cMvPolynomialEquiv, CMvPolynomial.C, CMvPolynomial.equiv_single, CMvPolynomial.equiv_C, CMvPolynomial.equiv_X}

\leanok

Use a direct sum of coefficients indexed by finitely supported natural exponent vectors,
represented with computable dependent finitely supported functions. There are executable
constants and variables, and an algebra equivalence to Mathlib's multivariate polynomials.
An unsafe representation instance extracts quotient representatives for display only;
it is not used by the algebraic algorithms or correctness proofs.

\end{definition}

\begin{lemma}[Normalization]

\label{lem:normalize}

\lean{CMvPolynomial.normalize, CMvPolynomial.normalize_eq}

\leanok

\uses{def:computable}

Rebuilding a polynomial from its nonzero coefficients removes accumulated zero terms
from its support bound and does not change its value.

\end{lemma}

\begin{proof}

\leanok

The linked Lean declarations contain the proof.

\end{proof}

\begin{definition}[Computable substitution and differentiation]

\label{def:compute-operations}

\lean{CMvPolynomial.bind, CMvPolynomial.bind_single, CMvPolynomial.equiv_bind, CMvPolynomial.setZero, CMvPolynomial.pderiv, CMvPolynomial.pderiv_single, CMvPolynomial.equiv_pderiv}

\leanok

\uses{def:computable}

Substitution is computed monomial by monomial and agrees under the algebra equivalence
with polynomial evaluation. Setting a variable to zero is a special substitution.
The derivative sends $rx^s$ to $r s_i x^{s-e_i}$, and agrees with formal partial
differentiation under the equivalence.

\end{definition}

\begin{definition}[Computable matrix polynomials]

\label{def:compute-matrices}

\lean{CMvPolynomial.rayleighDiff, CMvPolynomial.resultant, CMvPolynomial.charmatrix, CMvPolynomial.detPoly, CMvPolynomial.adjPoly, CMvPolynomial.sameFiber}

\leanok

\uses{def:compute-operations}

The computation model supplies Rayleigh differences, degree-one resultants,
$\operatorname{diag}(X)+A$, its determinant and adjugate, and a Boolean test for equality
of determinantal polynomials. Coefficients may be exact computable rings such as the
rationals. Matrix operations use Mathlib's determinant and adjugate algorithms.

\end{definition}

\begin{lemma}[Correctness of the computable operations]

\label{lem:compute-correct}

\lean{CMvPolynomial.equiv_setZero, CMvPolynomial.equiv_rayleighDiff, CMvPolynomial.equiv_charmatrix, CMvPolynomial.equiv_resultant, CMvPolynomial.equiv_detPoly, CMvPolynomial.equiv_adjPoly, CMvPolynomial.sameFiber_eq_true_iff_detPoly_eq}

\leanok

\uses{def:compute-matrices,def:setzero,def:resultant,def:adjpoly}

Transport through the algebra equivalence agrees with the mathematical specialization,
Rayleigh difference, characteristic matrix, resultant, determinantal polynomial and
polynomial adjugate. The Boolean comparison returns true exactly when $f_A=f_B$.
This certificate does not use the unfinished principal-minor coefficient theorem;
interpreting the Boolean test as principal-minor equality still uses
Theorem~\ref{thm:detpoly-expansion}.

\end{lemma}

\begin{proof}

\leanok

The linked Lean declarations contain the proof.

\end{proof}

\begin{definition}[Exact executable examples]

\label{def:compute-examples}

\lean{CMvPolynomial.Examples.A, CMvPolynomial.Examples.B, CMvPolynomial.Examples.D, CMvPolynomial.Examples.cutA, CMvPolynomial.Examples.cutB}

\leanok

\uses{lem:compute-correct,def:cut-swap}

Two rational two-by-two matrices related by diagonal conjugation have equal determinantal
polynomials; a third with a different off-diagonal product does not. Guards verify the
polynomial $x_0x_1+4x_0+x_1-2$, its derivative, specialization, Rayleigh difference,
a resultant and an adjugate entry, and normalization. A four-by-four pair formed by a
cut transpose has equal determinantal polynomials. Further guards check the empty
determinant and differentiation of $x_0^3$. The five matrices are exported definitions.

\end{definition}

\section{Problem-set results and alternative relations}

\begin{definition}[Independent row and column scaling]

\label{def:two-sided}

\lean{Matrix.TwoSidedDiagEquiv, Matrix.TwoSidedDiagEquiv.equivalence, Matrix.TwoSidedDiagEquiv.refl, Matrix.TwoSidedDiagEquiv.symm, Matrix.TwoSidedDiagEquiv.trans, Matrix.twoSidedDiagEquivSetoid, Matrix.TwoSidedDiagEquiv.transpose, Matrix.TwoSidedDiagEquivUpToTranspose, Matrix.TwoSidedDiagEquivUpToTranspose.equivalence, Matrix.TwoSidedDiagEquivUpToTranspose.refl, Matrix.TwoSidedDiagEquivUpToTranspose.symm, Matrix.TwoSidedDiagEquivUpToTranspose.trans, Matrix.TwoSidedDiagEquiv.twoSidedDiagEquivUpToTranspose, Matrix.twoSidedDiagEquivUpToTransposeSetoid}

\leanok

The alternative relation is
\[ M=\operatorname{diag}(l)N\operatorname{diag}(r) \]
for
independent families of units $l,r$, allowing rectangular matrices. It is an equivalence
relation and is preserved by transpose. For square matrices its closure under transpose
is also an equivalence relation. Both have the reflexive, symmetric, transitive and
setoid APIs, and the first relation embeds in the second.
This is a different relation from diagonal conjugacy in the principal-minor theorem.

\end{definition}

\begin{lemma}[Entries of diagonal conjugation]

\label{lem:diag-conj-entry}

\lean{diag_conj_apply}

\leanok

For a vector $d$ over a field, the entry formula is
\[ (\operatorname{diag}(d)A\operatorname{diag}(d^{-1}))_{ij}=d_iA_{ij}d_j^{-1}. \]
The entry formula itself does not need nonzero entries in $d$.

\end{lemma}

\begin{proof}

\leanok

The linked Lean declarations contain the proof.

\end{proof}

\begin{lemma}[Repeated rows, adjugates and inverses]

\label{lem:bach}

\lean{BachWarmup.det_updateRow_eq_zero, BachWarmup.adj, BachWarmup.adj_eq_adjugate, BachWarmup.mul_adj, BachWarmup.exists_inverse_iff_isUnit_det}

\leanok

Replacing one row by a distinct existing row gives determinant zero.
For indices $\{0,\ldots,n\}$, the cofactor definition
$\operatorname{adj}(A)_{ij}=(-1)^{i+j}\det A[\widehat{j},\widehat{i}]$ agrees with Mathlib's
adjugate and satisfies $A\operatorname{adj}(A)=\det(A)I$.
Over a commutative ring, a square matrix has a two-sided inverse exactly when its
determinant is a unit.

\end{lemma}

\begin{proof}

\leanok

Use determinant alternation for equal rows, the cofactor formula for the adjugate,
and the adjugate inverse when the determinant is a unit.

\end{proof}

\begin{definition}[Principal minors and mixed rows]

\label{def:diagonal-expansion}

\lean{DetDiagonal.principalMinor, DetDiagonal.mixRows}

\leanok

This is a problem-set proof of the expansion; the library version is
Theorem~\ref{thm:detpoly-expansion}.
In its separate namespace, principal minors are determinants of finite principal
submatrices. The matrix $\operatorname{mixRows}(S,M,N)$ uses rows of $M$ in $S$ and
rows of $N$ outside $S$.

\end{definition}

\begin{lemma}[Determinants of mixed rows]

\label{lem:mixed-rows}

\lean{DetDiagonal.row_diagonal_eq_smul_row_one, DetDiagonal.det_of_rows, DetDiagonal.det_mixRows_one, DetDiagonal.det_mixRows_diagonal}

\leanok

\uses{def:diagonal-expansion}

A diagonal row is a scalar multiple of the corresponding identity row. The determinant
is the alternating multilinear map on its row family. Mixing identity rows in $S$ with
rows of $A$ gives $A_{S^c}$. Mixing diagonal rows instead gives
$A_{S^c}\prod_{i\in S}x_i$.

\end{lemma}

\begin{proof}

\leanok

The linked Lean declarations contain the proof.

\end{proof}

\begin{theorem}[Determinant after adding a diagonal]

\label{thm:diagonal-expansion}

\lean{DetDiagonal.det_add_diagonal_formula}

\leanok

\uses{lem:mixed-rows}

For arbitrary $x:I\to R$,
\[ \det(A+\operatorname{diag}(x))=\sum_{S\subseteq I}A_{S^c}\prod_{i\in S}x_i. \]
This is the problem-set version of \texttt{Matrix.det\_diagonal\_add}
(Theorem~\ref{thm:detpoly-expansion}).

\end{theorem}

\begin{proof}

\leanok

Expand the determinant multilinearly in the sum of the original and diagonal row
families. Index the choices by subsets $S$, then evaluate each mixed-row determinant.

\end{proof}

\begin{theorem}[Number of terms in the diagonal expansion]

\label{thm:diagonal-expansion-count}

\lean{DetCounting.countState, DetCounting.Acount, DetCounting.Bcount, DetCounting.Acount_zero, DetCounting.Bcount_zero, DetCounting.Acount_one, DetCounting.Bcount_one, DetCounting.Bcount_succ, DetCounting.Acount_succ, DetCounting.Acount_second_order_additive, DetCounting.Bcount_closed_step, DetCounting.Acount_first_order, DetCounting.factorialQuotientCount, DetCounting.factorialQuotientCount_succ, DetCounting.Acount_eq_factorialQuotientCount, DetCounting.Acount_closed_form, DetCounting.rhsTermCount, DetCounting.choose_mul_complement_factorial, DetCounting.rhsTermCount_cast_eq_factorialQuotientCount, DetCounting.Acount_eq_rhsTermCount, DetCounting.counting_argument_final_nat, DetCounting.counting_argument_final_factorial}

\leanok

\uses{thm:diagonal-expansion}

Let $A_n$ and $B_n$ be defined by the recurrences $A_n=A_{n-1}+B_n$ and
$B_n=A_{n-1}+(n-1)B_{n-1}$, which count the terms of $\det(A+\operatorname{diag}(x))$ and
of $\det(A+\operatorname{diag}(0,x_2,\ldots,x_n))$ under cofactor expansion. Then
\[ A_n=\sum_{k=0}^n\binom nk(n-k)!=\sum_{k=0}^n\frac{n!}{k!}. \]
The identification of $A_n$ with the size of the support of an actual polynomial is not
formalized.

\end{theorem}

\begin{proof}

\leanok

Derive the first-order recurrence $A_{n+1}=(n+1)A_n+1$ and compare it with the recurrences
satisfied by the two closed forms.

\end{proof}

\begin{theorem}[Generalized Laplace on a finite ordinal]

\label{thm:fin-laplace}

\lean{Matrix.card_compl_fin, Matrix.det_eq_sum_neg_one_pow_mul_det_mul_det_fin}

\leanok

\uses{thm:generalized-laplace}

For $I=\{0,\ldots,N-1\}$ and $|S|=k$, one has $|S^c|=N-k$ and the increasingly
indexed expansion
\[ \det A=\sum_{|T|=k}(-1)^{\sum_{i\in S}i+\sum_{j\in T}j}
 \det A[S,T]\det A[S^c,T^c]. \]
This is the paper's finite-ordinal version of the general library theorem.

\end{theorem}

\begin{proof}

\leanok

Specialize the linearly ordered formula to a finite ordinal, where position equals value.

\end{proof}

\begin{definition}[Display instance]

\label{def:display-instance}

\lean{instReprCMvPolynomialOfDecidableEq}

\leanok

\uses{def:computable}

The representation instance prints terms with nonzero coefficients and their exponents.

Term order is unspecified. This instance is unsafe and is used only for display.

\end{definition}

\section{Additional matrix identities}

The following APIs are retained in \texttt{LEAN\_Warmup}; the proof of Theorem 2.3
does not depend on them.

\begin{lemma}[General cofactors and Laplace expansion]

\label{lem:warmup-laplace}

\lean{Matrix.det_updateCol_single_self, Matrix.det_updateRow_single, Matrix.det_updateCol_single, Matrix.adjugate_apply_eq_det_submatrix, Matrix.det_eq_sum_mul_det_submatrix_row, Matrix.det_eq_sum_mul_det_submatrix_col, Matrix.adjugate_apply_eq_det_submatrix_swap, Matrix.det_eq_sum_mul_det_submatrix_row_swap, Matrix.det_eq_sum_mul_det_submatrix_col_swap}

\leanok

\uses{lem:laplace}

General row and column deletion formulas use equivalences from an option type and the
sign of the induced permutation. Expanding the determinant gives row and column Laplace
formulas, including specializations matching the remaining indices by a swap.

\end{lemma}

\section{Generalized Laplace expansion}

\begin{lemma}[Group determinant terms by image of a row set]

\label{lem:laplace-group}

\lean{Matrix.sum_sign_mul_prod_eq_det_mul_det, Matrix.sum_sign_mul_prod_eq_sign_mul_det_mul_det, Matrix.det_eq_sum_det_mul_det_row, Matrix.det_eq_sum_det_mul_det_col, Matrix.sign_mul_det_mul_det_eq_of_map_eq}

\leanok

The terms of the Leibniz determinant formula indexed by permutations preserving $S$
sum to $\det A[S,S]\det A[S^c,S^c]$. Terms sending $S$ onto $\sigma(S)$ instead give
\[ \operatorname{sgn}(\sigma)
 \det A[S,\sigma(S)]\det A[S^c,\sigma(S^c)]. \]
The signed product depends only on the image set. Summing over column subsets of size
$|S|$ gives the row expansion; transposition gives the column expansion.

\end{lemma}

\begin{proof}

\leanok

Zero the off-block entries to identify the preserving-permutation sum with a block
diagonal determinant. Group the full Leibniz sum by the image of $S$.

\end{proof}

\begin{definition}[Chosen permutations for complementary minors]

\label{def:laplace-perm}

\lean{Matrix.exists_perm_map_eq, Matrix.laplacePerm, Matrix.map_laplacePerm, Matrix.det_eq_sum_laplacePerm_row, Matrix.det_eq_sum_laplacePerm_col}

\leanok

\uses{lem:laplace-group}

Equal-cardinality subsets admit a permutation sending one to the other.
Choose such a permutation for each pair and obtain row and column Laplace expansions
with these choices. No linear order on $I$ is required.

\end{definition}

\begin{lemma}[Signs and increasing enumerations]

\label{lem:shuffle-sign}

\lean{Matrix.sign_eq_prod_prod_lt, Matrix.sign_eq_of_strictMonoOn, Matrix.prod_orderEmbOfFin, Matrix.card_filter_lt_orderEmbOfFin, Matrix.prod_prod_ite_lt_mul, Matrix.card_compl_eq_of_card_eq}

\leanok

On a linearly ordered finite type, the sign of a permutation is a product of comparison
signs. If it is increasing separately on a subset and its complement, only cross-block
comparisons remain. Products and positions can be computed using increasing enumerations.
The auxiliary cardinality and comparison-product identities supply the sign calculation.

\end{lemma}

\begin{proof}

\leanok

The linked Lean declarations contain the proof.

\end{proof}

\begin{definition}[Order shuffle]

\label{def:shuffle}

\lean{Matrix.orderShuffle, Matrix.orderShuffle_orderEmbOfFin, Matrix.orderShuffle_orderEmbOfFin_compl, Matrix.map_orderShuffle, Matrix.sign_orderShuffle}

\leanok

\uses{lem:shuffle-sign}

For $|S|=|T|$, the order shuffle sends the increasing enumeration of $S$ to that of $T$
and likewise for their complements. It sends $S$ onto $T$ and its sign is
$(-1)^{\sum_{a\in S}\operatorname{pos}(a)+\sum_{b\in T}\operatorname{pos}(b)}$,
where $\operatorname{pos}(a)=|\{b:b<a\}|$.

\end{definition}

\begin{theorem}[Generalized Laplace expansion (Theorem 3.1)]

\label{thm:generalized-laplace}

\lean{Matrix.det_eq_sum_neg_one_pow_mul_det_mul_det}

\leanok

\uses{lem:laplace-group,def:shuffle}

For increasingly indexed minors,
\[ \det A=\sum_{|T|=|S|}(-1)^{\sum_{a\in S}\operatorname{pos}(a)+
 \sum_{b\in T}\operatorname{pos}(b)}\det A[S,T]\det A[S^c,T^c]. \]

\end{theorem}

\begin{proof}

\leanok

Use the order shuffle in the arbitrary-permutation expansion and substitute its sign.

\end{proof}
